\section{Introdução}
\label{sec:introducao}

Com o avanço das redes de comunicação e a convergência entre voz, vídeo e dados
em uma única infraestrutura, surgiu a necessidade de mecanismos capazes de
garantir a qualidade na transmissão de informações. Historicamente, as redes de
negócios operavam como entidades separadas — chamadas telefônicas trafegavam em
redes dedicadas, enquanto dados percorriam infraestruturas distintas
\cite{fortinet2024}. Com o tempo, essa separação cedeu lugar à integração, e
aplicações interativas como videoconferência, \textit{Voice over IP} (VoIP) e
\textit{streaming} passaram a compartilhar a mesma infraestrutura de rede.

Esse cenário introduziu desafios consideráveis. Aplicações multimídia são
altamente sensíveis a atrasos, variações de atraso e perdas de pacotes, exigindo
que a rede seja capaz de priorizar certos tipos de tráfego em detrimento de
outros \cite{techtarget2024}. Além disso, o crescimento da Internet das Coisas
(IoT) ampliou ainda mais essa demanda, uma vez que sensores e dispositivos
industriais dependem de comunicação em tempo real para operar corretamente
\cite{fortinet2024}.

Nesse contexto, a Qualidade de Serviço (QoS) emerge como um conjunto de
mecanismos e tecnologias fundamentais para garantir a \textit{performance}
adequada de aplicações críticas, mesmo em ambientes de rede com recursos
limitados. Este trabalho apresenta os principais conceitos, elementos avaliados,
técnicas e aplicações do QoS em redes modernas.
